\documentclass[11pt]{article}
\usepackage{icmlko}
\begin{document}

\title{둘 다 아니면 아무것도 --- 검사를 통과한 도메인은 안 벌고, 실패한 도메인도 혼자서는 못 벌었다}
\author{월드모델 실험실}
\date{노트 565 · 566 · 2026년 8월 5일}
\maketitle

\begin{abstract}
노트 $563$ 이 ``게임의 검색 관측-$0$ 은 전부 한글 제목''을 찾다가 새 축이
따라 나왔다 --- \texttt{lang\_ko}(제목에 한글이 있나). 작품의 성질이고 출시
시점 확정이라 사전이며, 라벨을 한 번도 안 본다. \texttt{tagaxes} 채택 검사
셋을 \textbf{모바일에서 다 통과}했다(① 학습 $+0.4362$ · ② 다섯 조각
$[+,+,+,+,+]$ · ③ 최대 겹침 $0.371$).

\begin{center}
\begin{tabular}{lrr}
\hline
 & 짝 씨앗 $12$ & \textbf{짝 씨앗 $40$} \\
\hline
\textbf{모바일} & $+0.0021$ ($10/12$) & $\mathbf{+0.0007}$ ($23/40$, 문다) \\
\textbf{비게임 앱} & $-0.0015$ (문다) & $\mathbf{-0.0028}$ ($0$ 밖 음수) \\
\hline
\end{tabular}
\end{center}

\textbf{이득은 사라지고 손해는 굳었다.} 관문 ③(어떤 자도 $0$ 밖 음수가
아니다)을 못 넘고 노트 $426$ 거울도 미달이다 --- \textbf{기각}.
\end{abstract}

\begin{figure}[t]
\centering
\includegraphics[width=\textwidth]{figs/f1.pdf}
\caption{\textbf{가} --- 씨앗을 늘리면. \textbf{나} --- 둘 다 아니면 아무것도.
\textbf{다} --- 검사의 한계.}
\end{figure}

\section{둘 다 아니면 아무것도}

\begin{center}
\begin{tabular}{lrr}
\hline
축을 주는 곳 & 판 & KR 만화 \\
\hline
모바일 $\cdot$ 게임 둘 다 & $+0.0006$ (문다) & $\mathbf{+0.0053}$ ($38/40$) \\
\textbf{모바일에만} & $\mathbf{-0.0055}$ ($0/12$) & $+0.0016$ \\
\textbf{게임에만} & $\mathbf{-0.0049}$ ($0/40$) & $+0.0011$ \\
\hline
\end{tabular}
\end{center}

\textbf{한쪽만 주면 판이 $-0.005$ 떨어진다.} 그 축이 \emph{그 도메인에서만}
변하면 --- 나머지는 값이 상수이므로 --- \textbf{도메인 지시자}가 된다.
노트 $540$ 이 ``한 도메인에만 있는 열은 그 도메인의 지시자다''로 닫은 것이
\textbf{결측 무늬가 아니라 값 열에서도} 성립한다.

\begin{quote}
그러면 ``둘 다''의 이득($+0.0053$)이 무엇인지 남는다. \textbf{게임 몫이
아니다}(게임에만 주면 $+0.0011$). \textbf{모바일 몫도 아니다}(모바일에만
$+0.0016$, 그리고 모바일 자신은 $+0.0007$ 로 안 오른다). 남는 해석은
\textbf{``두 도메인에서 변해야 지시자가 아니게 된다''} 뿐이고, 그것은 정보가
아니라 \emph{배선}이다.
\end{quote}

\section{채택 검사는 어디서 벌지 못 맞힌다}

\begin{center}
\begin{tabular}{lll}
\hline
 & 검사 & 결과 \\
\hline
\textbf{모바일} & ① $+0.4362$ · ② $[+,+,+,+,+]$ · ③ $0.371$ & \textbf{셋 다 통과} \\
게임 & ① $-0.0489$ · ② $[+,+,-,-,-]$ & ①② 실패 \\
\hline
\end{tabular}
\end{center}

검사를 \textbf{통과한} 모바일이 씨앗 $40$ 에서 $+0.0007$(문다)이고, 검사를
\textbf{실패한} 게임은 혼자 주면 판을 $-0.0049$ 떨어뜨린다.

\begin{quote}
\textbf{검사 ①②③ 은 후보를 거르는 데는 쓸모 있지만 \emph{어디서 벌지}는
못 맞힌다.} 그것은 ``이 축이 이 도메인에서 라벨과 붙는가''를 묻는데,
실제로 판이 움직이는 것은 \emph{공유 모형 안에서 그 열이 어떻게 쓰이는가}로
정해진다.
\end{quote}

\section{씨앗 $12$ 로 채택했으면 틀렸다}

\begin{itemize}
\item 모바일 $+0.0021$($10/12$, $0$ 밖 양수) $\to$ $+0.0007$($23/40$, 문다)
--- \textbf{이득이 사라졌다}
\item 비게임 앱 $-0.0015$(문다) $\to$ $-0.0028$($0$ 밖 음수) ---
\textbf{손해가 굳었다}
\end{itemize}

노트 $518$(어느 자에서든 결정이 $0$ 근처면 씨앗 $40$)이 이번엔 \emph{막는}
쪽으로 값을 했다. 바로 앞 노트 $553 \cdot 564$ 에서는 씨앗 $40$ 이
\emph{통과시키는} 쪽으로 값을 했다 --- \textbf{같은 조항이 양쪽으로 작동한다.}

\section*{그래도 남는 것}

기각이지만 \texttt{lang\_ko} 가 헛되지는 않았다. 게임에서 그 축은
\texttt{trend\_level} 과 \textbf{$+0.519$}(관측 행 한정)로 겹친다 --- 즉
\textbf{검색 값의 흐린 사본}이다. 그런데 노트 $564$ 가 게임의 관측-$0$
$42$ 행을 결측으로 돌렸으므로, \textbf{그 $42$ 행에서만 검색이 침묵하고
\texttt{lang\_ko} 가 말한다}. 이득의 자리가 거기였고, 그것으로는
$0.005$ 를 못 넘었다.

\section*{조항}

\begin{quote}
\textbf{새 축이 한 도메인에서만 변하면 그것만으로 해롭다} --- 노트 $540$ 의
도메인 지시자가 \emph{값 열에서도} 성립한다. \textbf{그리고 이득이 ``둘 다
있을 때만'' 난다면 정보가 아니라 배선을 보고 있는 것이다} --- 단일 도메인
팔로 갈라야 안다.
\end{quote}

\section*{남은 것}

\begin{itemize}
\item \textbf{관측-$0$ 이 언어 하위집단인 도메인은 게임 하나뿐}이다 ---
값있음 대 관측-$0$ 의 한글 비중 차가 게임 $+64$\%p 이고 나머지는
$0$\%p(도메인 전체가 한글이거나 전체가 비한글) 또는 모바일 $-13$\%p 다.
\textbf{노트 $439$ 의 남은 것을 닫는다.}
\item \textbf{수집 순서는 코드로 끊었다} --- \texttt{ingest/trend\_all.collect}
가 씨앗 고정 섞기 뒤에 자른다. 노트 $379$(펀딩) $\cdot$ $554$(만화)가 같은
결함을 두 번 겪었고 \textbf{세계애니 파일도 완전 정렬}이라 세 번째가 예약돼
있었다.
\item \textbf{이웃 효과를 잡는 손잡이는 여전히 없다} --- 노트 $562$ 가
모형 고르기를 닫았고(검증이 다섯 배 부족), 여기서 새 축도 닫혔다.
\end{itemize}

\end{document}
